% V30 independent methods technical note; RETROSPECTIVE, not submitted.
\documentclass[11pt,a4paper]{article}
\usepackage[margin=0.92in]{geometry}
\usepackage[T1]{fontenc}
\usepackage{amsmath,amssymb,booktabs,graphicx}
\usepackage[hidelinks]{hyperref}
\usepackage{lmodern}
\usepackage[protrusion=true,expansion=false]{microtype}
\title{EEG Recording Identity Beyond Subject Identifiers:
Exact Epoch Overlap and Matched Model--Selection Audits}
\author{AUTHOR IDENTITIES AND AFFILIATIONS TO BE VERIFIED}
\date{}
\begin{document}\maketitle
\begin{abstract}
\textbf{Objective.} Subject-disjoint EEG validation assumes that released subject
identifiers correspond to independent recordings. We test this assumption in
a public acute-stroke motor-imagery acquisition and quantify how record-content
purging interacts with source-model fitting and regularisation choice.
\textbf{Methods.} Under pinned MOABB 1.7.2 processing, exact cross-ID EEG
epoch/label fingerprints define record-content groups. All recording groups
are checked in three public acquisitions. In the affected cohort, we
reproduce original and fully purged held-record AUC values, then evaluate
a frozen $2\times2$ model-fit by shrinkage-selection factorial on the same
49 held recording groups.
\textbf{Results.} Liu2024 contains 50 released IDs and 2,000 processed trial
instances, but 1,960 unique exact trial contents; IDs 4 and 5 share all
40 epochs in the fixed pipeline. No cross-ID exact overlaps were observed
in Weibo2014 (10 IDs; 1,580 trials) or BNCI2014\_002 (14 IDs; 2,240 trials).
Source-macro held-record AUC changed from 0.594000 (nominal 50-ID average)
to 0.588673 (49 identical held targets with original fits) and 0.582551
(49-source-group full refit). Among matched 49 groups, the source-macro
regularisation choice changed in 37, with opposing average conditional
selector effects ($-0.001327$ versus $+0.001786$) on the two fitted models.
The correctly predictive source-mean Gaussian policy changed its selected
shrinkage in zero groups.
\textbf{Conclusion.} Subject IDs alone may not ensure recording-content
independence. The original-to-corrected AUC difference mixes a reporting
denominator shift, model changes and selection interactions; it is neither
a patient-level causal leakage penalty nor a new transfer gain.
\end{abstract}

\noindent\textbf{Keywords:} EEG; motor imagery; recording identity;
cross-subject validation; exact epoch overlap; regularisation.

\section{Introduction}
Cross-subject EEG generalisation requires a defensible definition of
independent samples. Recent accepted work demonstrates large optimism
under window-level rather than subject-level validation and proposes
hierarchical statistical evaluation \cite{gil2026subjectvalidation}.
Even when a dataset provides subject IDs, the stronger claim that
different IDs never reference overlapping recording contents should be
verified when the scientific estimand is an unseen-recording result.

The public Liu2024 motor-imagery acquisition is described as comprising
50 acute-stroke patients \cite{liu2024stroke}. This study audits numerical
EEG trial contents available through a reproducible public software
pipeline; it neither identifies clinical participants nor disputes the
underlying recruitment metadata. The research questions are whether
exact epoch contents recur under different released IDs and whether
their removal affects decoding through altered evaluation weights,
source fitting, source regularisation choice, or interactions.

\section{Methods}
\subsection{Preprocessing, exact identity and content-safe splits}
We freeze MOABB 1.7.2, binary left/right imagery, original 29 EEG
channels, 8--30 Hz filtering and 160 Hz resampling for Liu2024.
Each of the 50 IDs contributes 40 processed trials. Exact numeric
trial content and event label are compared across the whole
public processed cohort; an edge joins IDs with matching epochs.
Connected components are treated as recording-content groups. If two
released IDs contain all 40 matching trials, those ID groups cannot
appear on opposite sides of any claimed recording-disjoint split.

This is an audit of processed arrays, not evidence of identity
of actual patients or raw EDF/MAT data. Global feature/epoch
inspection to discover the groups means these retrospective
experiments are not strict physical target-isolated Z0 studies.

Two already-exposed cohorts were fixed before running identical
cross-ID fingerprint controls: Weibo2014 (C3,Cz,C4, left/right hand)
and BNCI2014\_002 (15 EEG channels, right hand/feet). Their inclusion
tests whether the finding appears uniformly across different acquisitions.
It does not supply new decoder-generalisation confirmation.

\subsection{Matched full-refit and source-selector cross}
The original nominal-50-ID source model excludes one target ID, while
possibly retaining a same-content counterpart. The purged
49-recording-group model excludes the entire held content group and
refits its source covariance tangent reference, spectral LDA and source
selection on the remaining 48 representative recordings.
All fitted-model comparisons evaluate the SAME 49 held-recording
targets. Both parameter choices are copied from separately frozen
source-only artifacts, not optimised with held target AUC.

For an original fit $F_A$, purged fit $F_B$, old source
choice $\lambda_A$ and purged choice $\lambda_B$, define
\[
\begin{array}{c|cc}
 & \lambda_A & \lambda_B\\\hline
 F_A&Y_{AA}&Y_{AB}\\
 F_B&Y_{BA}&Y_{BB}
\end{array}.
\]
The software-factorial identities are
\[
Y_{BB}-Y_{AA}=(Y_{AB}-Y_{AA})+(Y_{BB}-Y_{AB})
=(Y_{BA}-Y_{AA})+(Y_{BB}-Y_{BA}),
\]
and the model--selector interaction is
$I=Y_{BB}-Y_{BA}-Y_{AB}+Y_{AA}$.
These are arithmetic decompositions, not causal identification.
Every one of the $49\times3$ held-group AUC comparisons at
$Y_{AA}$ and $Y_{BB}$ is independently checked against
two earlier immutable outcome artifacts to $10^{-10}$.
The inferential unit for descriptive reports is one processed
recording group, not a trial or confirmed independent patient;
LOPO source training pools overlap heavily.

\section{Results}
\subsection{Identity across three EEG acquisitions}
\begin{table}[ht]\centering\small
\caption{Exact processed trial-content audit under preselected
pipelines. Unique contents do not imply uniquely identified patients.}
\begin{tabular}{lrrrr}
\toprule
Acquisition & Released IDs & Trial instances & Unique & Cross-ID edges\\\midrule
Liu2024 &50&2,000&1,960&1\\
Weibo2014 &10&1,580&1,580&0\\
BNCI2014\_002 &14&2,240&2,240&0\\
\bottomrule
\end{tabular}\end{table}
In the processed Liu2024 signals, IDs 4 and 5 have the same complete
40-trial EEG and label contents. No further partial cross-ID exact
trial repetitions were found among the 2,000 trial instances.
Zero exact overlaps in the two controls cannot rule out near-duplicates
or overlapping source windows with different sampled arrays.

\subsection{Matched fit--parameter factorial and reporting denominator}
The original source-macro average of 0.594000 over 50 nominal IDs
declined to 0.588673 after retaining only 49 recording representatives
with ORIGINAL fits and lambdas, a denominator-only contrast of
$-0.005327$. Under a fully purged, newly fitted source model,
the matched 49-record mean was 0.582551, an additional
$-0.006122$. The aggregate $-0.011449$ difference is
not by itself a causal EEG leakage-effect estimate.

\subsection{Real 49-record factorial diagnosis: denominator versus refitting versus selection}
\label{sec:v30-actual-crossed}
The previously locked V29 experiment completed all 49 unique
recording-group held-target evaluations with the originally
frozen model and parameter choices, and independently
refitted all 49 content-disjoint models.
For each group, both frozen sets of selected shrinkage
values were evaluated under both fitted source models.
The original per-record target AUC values and the purged
per-record target AUC values were checked to numerical
tolerance $10^{-10}$ against their independent immutable
archives; all $49\times3$ method comparisons passed.
This is retrospective, previously exposed EEG, not a
new unseen-target performance experiment.

\begin{table}[tb]
\centering
\small
\caption{Liu2024 source macro-AUC shrinkage policy: genuine
paired $2\times2$ model fit $\times$ frozen parameter analysis
on the same 49 held recording groups. $A$ is the original
nominal-50-ID source fit, $B$ the fully record-disjoint 49-group
refit. The reported values are held-record macro-AUC,
not independent trials or new patients.}
\label{tab:v30-crossed}
\begin{tabular}{lcc}
\toprule
Model fit & Original choice $\lambda_A$ & Purged choice $\lambda_B$\\
\midrule
Original $F_A$ & 0.588673 & 0.587347\\
Purged $F_B$ & 0.580765 & 0.582551\\
\bottomrule
\end{tabular}
\end{table}

The source-macro selector changed its frozen $\lambda$
in 37/49 groups. Its mean \emph{conditional parameter
substitution} was $-0.001327$ AUC on the original
fit but $+0.001786$ on the purged fit, producing
an interaction $+0.003112$. The \emph{total matched}
change from old fit/old choice to purged fit/purged
choice was $-0.006122$. In contrast, the correct
source-trained-mean predictive Gaussian and historical
validation-class-centred covariance policies selected
the same shrinkage under both group specifications
in all 49 cases. Their corresponding model-refit-only
matched differences were $-0.003724$ and $-0.006071$.
Thus the observed change should not be labeled a
universal selector instability: the source-macro
choice reacts to record composition, but the
predictive choice did not react in this test.

\paragraph{Do not confuse denominators with a paired treatment effect.}
The original, contaminated \emph{50-ID} mean
source-macro AUC was 0.594000, whereas the subset of
49 recording representatives under the unchanged
original fit and unchanged parameter choices
averaged 0.588673. The $0.005327$ difference is
a \emph{reporting-denominator contrast} caused by
removing one duplicated held-test identifier.
The additional matched $0.006122$ difference to
the properly purged and refitted result yields
the descriptive full $0.011449$ difference,
but each term is contingent on the original
dataset and model. It is incorrect to report
$0.011449$ as a causal leakage bias, a
cross-subject generalization gain, or a
statistically independent 50-versus-49 test.

All 49 held-record $Y_{AA}$ and $Y_{BB}$ values
were independently reproduced, not just their
means. The completed workflow is GitHub Actions
\texttt{38290349868}, receipt SHA-256 prefix
\texttt{0e7bfdb73970fcff}. This is an
observable software-factorial decomposition,
not an intervention identifying the direct
causal contribution of duplicated EEG signals.


\section{Discussion}
Our observation is more specific than generic random-window leakage:
the public released label unit was not enough to ensure disjoint
processed signal contents. The full data-integrity audit and two
independent-acquisition controls make the particular duplicate
verifiable and falsifiable, while the same-held-target factorial
prevents silently attributing the whole correction to a changing
regularisation parameter.

The source-macro policy changed $\lambda$ in 37/49 groups,
yet its two conditional source-model selector effects had
opposite signs. In contrast, neither the source-trained-mean
predictive Gaussian nor historical class-centred covariance
choice changed its $\lambda$ under this comparison. A changed
target mean could thus arise even with \emph{no}
regularisation-selection changes. These results separate
source-model dependence from parameter substitution, but do
not uniquely isolate causes of training contamination.

The overlap finding pertains to a fixed processed software
representation. Authoritative resolution of whether patient
records themselves were duplicated requires the original
data custodians and cannot be inferred from processed EEG.
All performance outcomes had been exposed before the V29
factorial and all apparent group-bootstrap intervals are
descriptive, not new confirmatory tests. This analysis does
not demonstrate a SOTA EEG decoder, a clinical outcome, or
a safety guarantee for unseen stroke patients. The
original Default Regularisation EEG estimand analysis
remains a separate primary paper and should not combine
these findings into its earlier frozen Holm families.

\section{Conclusion}
An EEG subject label is not automatically an independent
signal identity. Full processed epoch-content checks found
one complete cross-ID duplicate in a 50-ID acute-stroke
motor-imagery acquisition, and no exact cross-ID matches
in two preselected other acquisitions. A 49-recording-group
full refit and $2\times2$ frozen fit--selector contrast
showed that reporting weights, model fitting and source
selection must not be collapsed into a single asserted
leakage penalty. The practical implication is to audit
recording provenance before interpreting cross-subject
regularisation comparisons.

\section*{Statements and Declarations}
\textbf{Data:} Liu2024, Weibo2014 and BNCI2014\_002
are publicly distributed; code and immutable numerical
experiment receipts are provided in the public GitHub
repository. No EEG signals are included in this manuscript.\\
\textbf{Author, ethics, competing interests,
funding and AI-use statements:} to be verified by
the actual authors before submission; no claim is made
on anyone's behalf.

\bibliographystyle{unsrt}
\bibliography{references_v10}
\end{document}
